\section{From Constitution to Missal}

The Council closed on 8 December 1965. The Missal appeared in 1970. Everything in between was the work not of the Council but of the Apostolic See exercising the authority article 22 reserved to it---a distinction with consequences, because praise and blame directed at ``the Council's Missal'' land, on inspection, on a series of papal and curial acts. The revision was conducted by a commission Paul VI established in January 1964 to implement the constitution (the \emph{Consilium ad exsequendam Constitutionem de Sacra Liturgia}), working through drafting groups of liturgical scholars, with instalments of reform issued across the decade: the vernacular's first expansions, the simplifications of 1964 and 1967, new Eucharistic Prayers in 1968.\endnote{The Consilium was established by the motu proprio \emph{Sacram Liturgiam} (25 January 1964); the staged instruments include \emph{Inter oecumenici} (1964) and \emph{Tres abhinc annos} (1967). This chronology is orienting background at the level of common documentary record; the argument rests on the promulgating acts quoted below, and no claim here depends on the staged instruments' exact content.} The culminating acts are two, and both deserve exact quotation because each is routinely misdescribed.

The first is Paul VI's apostolic constitution \emph{Missale Romanum} of 3 April 1969. It opens by honoring what it revises---the 1570 Missal of Pius V, fruit of Trent, which ``innumerable holy men'' had used to nourish their piety---and then presents the new book explicitly as the same kind of act: as Pius V promulgated an ``instrument of liturgical unity'' by decree of Trent, so Paul VI promulgates a Missal renewed by decree of Vatican II.\endnote{Paul VI, apostolic constitution \emph{Missale Romanum} (3 April 1969), official English web text checked 2026-07-25; Latin text as reprinted in the 2002 third typical edition checked the same day at the registered artifact (physical pages 4--7). The constitution's opening account of the conciliar basis quotes \emph{Sacrosanctum Concilium} 21, 50, 51, and 58 (its notes 4--7), and its lectionary paragraph quotes article 51 again.} Its own summary of the ``new composition'' names four things: a General Instruction prefacing the book with ``the new regulations \ldots{} for the celebration of the Eucharistic Sacrifice''; the Eucharistic Prayer enriched by a great number of Prefaces and by three added Canons---``we have decided to add three new Canons to this Prayer''---with the words of the Lord made identical in every anaphora; the rite of Mass simplified, ``due care being taken to preserve their substance''; and the readings reorganized so that, per article 51 of the constitution, ``a more representative portion of the Holy Scriptures'' is read, on Sundays in a three-year cycle with an Old Testament reading added before Epistle and Gospel.\endnote{\emph{Missale Romanum} (1969), as above, checked 2026-07-25. The prescribed common form of the words over the bread and chalice is set out in capitals in the constitution's text (\latin{Accipite et manducate ex hoc omnes\ldots{}; Accipite et bibite ex eo omnes\ldots}). The three added Eucharistic Prayers stand beside the Roman Canon, which remains as Eucharistic Prayer I; the constitution's Latin force clause (\latin{Nostra haec autem statuta et praescripta nunc et in posterum firma et efficacia esse et fore volumus}) was checked against the 2002 reprint.} Its hope is stated in the vocabulary of unity, not experiment: that the Missal ``will be received by the faithful as an instrument which bears witness to and which affirms the common unity of all,'' so that ``in the great diversity of languages, one unique prayer will rise'' to the Father---a sentence which shows that by 1969 the pope himself, not a runaway bureau, envisioned habitual vernacular celebration. The prescriptions were ordered into effect from 30 November 1969, the First Sunday of Advent.

The second act is the decree by which the Sacred Congregation for Divine Worship, ``by mandate of the Supreme Pontiff,'' promulgated the finished book on 26 March 1970 and declared it the typical edition: \latin{novam hanc editionem Missalis Romani ad normam decretorum Concilii Vaticani II confectam promulgat et uti typicam declarat}---promulgates this new edition of the Roman Missal, prepared according to the norms of the decrees of the Second Vatican Council, and declares it typical.\endnote{Sacred Congregation for Divine Worship, decree \emph{Celebrationis eucharisticae} (Prot.\ n.~166/70, 26 March 1970), Latin text as reprinted in the 2002 third typical edition, checked 2026-07-25 at the registered artifact (physical page 1). The decree permits the Latin edition's use upon publication and commits vernacular editions and their dates to the episcopal conferences with Apostolic See confirmation.} The edition history after 1970 is short and matters for exactness: an emended reprint in 1971; a second typical edition declared on 27 March 1975, its changes catalogued in the promulgating act (ministries of acolyte and lector replacing the subdiaconate in the Institutio, added ritual and votive formularies); the third typical edition approved by John Paul II on 10 April 2000 and declared on 20 April 2000, published in 2002; and an emended reprint of the third edition in 2008---a reprint, the official notice is careful to say, not a new edition.\endnote{The edition-history line (\latin{Editio typica, 1970; Reimpressio emendata, 1971; Editio typica secunda, 1975; Editio typica tertia, 2002}) and the promulgating acts of 1970, 1975, and 2000 were checked 2026-07-25 in the prefatory matter of the registered 2002 reproduction (physical pages 1--3); the 1975 act is Prot.\ N.~1970/74 (27 March 1975), the 2000 decree Prot.\ N.~143/00/L. For 2008, the Congregation's notice in \emph{Notitiae} 44 (2008) 367 identifies the \latin{reimpressio emendata} and its bounded corrections (registered with the \latin{variationes} list at pp.~368--387); no claim is made beyond that notice. Vernacular missals are territorial books confirmed by the Apostolic See and are not otherwise treated.}

Set the finished Missal beside the constitution, and the honest result is a column of each. Squarely mandated: the revised Order of Mass with simplified rites (art.~50), the enlarged lectionary (art.~51), concelebration and Communion under both kinds within limits (arts.~55, 57), the homily and common prayer (arts.~52--53). Authorized by the constitution's mechanisms though not commanded by its text: vernacular celebration in the people's parts and eventually throughout, which entered by territorial decisions confirmed by Rome under articles 36 and 54 and stood presumed in the 1969 constitution's ``great diversity of languages.'' Nowhere in the conciliar text, and resting on distinct postconciliar instruments: celebration facing the people, Communion in the hand, and the effective disappearance in ordinary parishes of the Latin and chant that articles 36, 54, and 116 expected to persist.\endnote{For the bounded negative claims about the constitution's text, see the note in section 3. Communion in the hand entered by the instruction \emph{Memoriale Domini} (29 May 1969), identified here by title and date only. Celebration facing the people is a matter of church arrangement and practice, not a rubrical obligation of the Missal; its history is outside this essay's scope.} And in one respect the reform confessedly went beyond restoration to composition: the three added Eucharistic Prayers stood beside a Roman Canon that had stood textually fixed, in Paul VI's own words, since the fourth and fifth centuries.

That sorting, dull as it is, dissolves two rhetorically convenient identities at once. ``The Mass of the Council'' is not simply the Council's: the Fathers voted principles and limits; a papal commission, then a pope, made the determinations, many of which---the scale of new composition, the pace, the effective retirement of the previous book from parish life---cannot be read off the conciliar text and must be defended, if defended, as prudent exercises of the authority article 22 confers. But neither is the Missal a rogue product: every operative act from 1963 to 1970 was performed by the Apostolic See, the one authority the constitution names as competent, and the book's structure answers point for point to mandates genuinely in the text. Whoever argues that the reform betrayed the Council argues against the pope's judgment, not against a forger; whoever argues that the Council commands every feature of postconciliar practice must produce articles that do not exist. What neither sorting settles is whether any of this---mandated, authorized, or added---explains what happened to Catholic life in the decades that followed. For that, the numbers must be faced.
